\documentclass[conference]{IEEEtran}

% Pacotes basicos
\usepackage[utf8]{inputenc}
\usepackage[T1]{fontenc}
\usepackage[brazil]{babel}
\usepackage{cite}
\usepackage{amsmath,amssymb,amsfonts}
\usepackage{graphicx}
\usepackage{url}
\usepackage{booktabs}
\usepackage{array}
\usepackage{listings}
\usepackage{xcolor}

% Configuracao para exibicao de codigo
\lstdefinestyle{pythonstyle}{
    language=Python,
    basicstyle=\ttfamily\footnotesize,
    breaklines=true,
    frame=single,
    numbers=left,
    numberstyle=\tiny,
    captionpos=b,
    showstringspaces=false
}

\lstdefinestyle{asmstyle}{
    basicstyle=\ttfamily\footnotesize,
    breaklines=true,
    frame=single,
    numbers=left,
    numberstyle=\tiny,
    captionpos=b,
    showstringspaces=false
}

\renewcommand\IEEEkeywordsname{Palavras-chave}
\renewcommand{\lstlistingname}{Listagem}

\begin{document}

\title{Desenvolvimento de um Assembler para a Arquitetura RISC-V32I}

\author{
\IEEEauthorblockN{Demétrius Mendes Miranda Rocha}
\IEEEauthorblockA{232012974}
\and
\IEEEauthorblockN{Lucas Silva Nóbrega}
\IEEEauthorblockA{180035096}
\and
\IEEEauthorblockN{João Gabriel Melo de Lima}
\IEEEauthorblockA{241032617}
\and
\IEEEauthorblockN{Marcus Vinicius Balbino Cavalcante}
\IEEEauthorblockA{160135974}
\and
\IEEEauthorblockN{Samuel Carvalho Dos Santos}
\IEEEauthorblockA{211010388}
}

\maketitle

\begin{abstract}
Este artigo apresenta o desenvolvimento de um assembler para um subconjunto da arquitetura RISC-V32I. A aplicação foi implementada em Python, em arquivo único, com o objetivo de processar arquivos Assembly com extensão \texttt{.asm}, interpretar instruções e dados, resolver rótulos e gerar arquivos de inicialização de memória no formato MIF. O trabalho aborda os fundamentos da codificação de instruções, a organização das seções \texttt{.text} e \texttt{.data}, a metodologia adotada na implementação e os resultados obtidos com a geração dos arquivos \texttt{.text.mif} e \texttt{.data.mif}. Dessa forma, o projeto relaciona conceitos de Organização e Arquitetura de Computadores, linguagem Assembly, representação hexadecimal, manipulação de arquivos e funcionamento básico de um assembler.
\end{abstract}

\begin{IEEEkeywords}
RISC-V32I, Assembly, assembler, Python, MIF, arquitetura de computadores.
\end{IEEEkeywords}

\begin{table}[!ht]
\centering
\caption{Informa\c{c}\~{o}es gerais do laborat\'{o}rio}
\label{tab:identificacao}
\begin{tabular}{>{\bfseries}p{0.34\linewidth} p{0.52\linewidth}}
\toprule
Disciplina & Organiza\c{c}\~{a}o e Arquitetura de Computadores \\
Turma & 03 \\
Professor & Fl\'{a}vio Vidal \\
Semestre & 1\textordmasculine/2026 \\
Laborat\'{o}rio & Laborat\'{o}rio 01 \\
Tema & Assembler para Assembly RISC V32I \\
\bottomrule
\end{tabular}
\end{table}

\section{Introdução}

A linguagem Assembly é uma representação simbólica de baixo nível utilizada para descrever instruções executadas diretamente por um processador. Por permitir uma relação mais direta com registradores, memória e fluxo de execução, seu estudo é fundamental em Organização e Arquitetura de Computadores.

Na arquitetura RISC-V, as instruções são codificadas em palavras de 32 bits e seguem formatos padronizados, como R, I, S, B, U e J. A variante RV32I define um conjunto básico de instruções inteiras, incluindo operações aritméticas, lógicas, acessos à memória, desvios condicionais e saltos. Cada formato determina a posição dos campos necessários para a codificação, como \texttt{opcode}, \texttt{rd}, \texttt{rs1}, \texttt{rs2}, \texttt{funct3}, \texttt{funct7} e imediato.

O processo de montagem, realizado por um assembler, consiste em converter instruções escritas em Assembly para código de máquina. Para isso, é necessário identificar mnemônicos, interpretar operandos, reconhecer registradores, tratar valores imediatos, processar rótulos e organizar as seções do arquivo, principalmente \texttt{.text}, referente às instruções, e \texttt{.data}, referente aos dados.

Neste projeto, foi desenvolvido um assembler em Python para um subconjunto da arquitetura RISC-V32I. A aplicação recebe um arquivo Assembly com extensão \texttt{.asm}, processa as instruções e os dados declarados e gera arquivos de saída no formato \texttt{.mif}, utilizado para inicialização de memórias em ambientes de simulação e projetos digitais.

Dessa forma, o projeto aplica conceitos como codificação de instruções, organização de memória, representação hexadecimal, manipulação de arquivos e funcionamento básico de um assembler, permitindo compreender as etapas necessárias para transformar um programa Assembly em uma representação compatível com o hardware ou com ambientes de simulação.

\section{Objetivos}

O objetivo principal deste laborat\'{o}rio \`{e} desenvolver uma aplica\c{c}\~{a}o capaz de atuar como assembler para um subconjunto de instru\c{c}\~{o}es da arquitetura RISC V32I. A ferramenta deve receber um arquivo de entrada em Assembly, processar as se\c{c}\~{o}es de texto e dados, traduzir as instru\c{c}\~{o}es para hexadecimal e gerar arquivos de sa\'{i}da no formato MIF.

De forma espec\'{i}fica, o projeto busca compreender a estrutura da arquitetura RISC V32I, interpretar instru\c{c}\~{o}es Assembly, reconhecer registradores, imediatos e r\'{o}tulos, montar os campos bin\'{a}rios das instru\c{c}\~{o}es suportadas, gerar a representa\c{c}\~{a}o hexadecimal correspondente e produzir arquivos de inicializa\c{c}\~{a}o de mem\'{o}ria para as regi\~{o}es de instru\c{c}\~{o}es e dados.


\section{Fundamentação Teórica}

\subsection{Arquitetura RISC-V32I}

A arquitetura RISC-V é baseada no princípio RISC, caracterizado por um conjunto de instruções reduzido, regular e organizado para simplificar a implementação do processador. Neste projeto, utiliza-se a variante RV32I, que corresponde ao conjunto base inteiro de 32 bits. Nessa arquitetura, as instruções possuem tamanho fixo de 32 bits e operam sobre 32 registradores inteiros de propósito geral, identificados de \texttt{x0} a \texttt{x31} \cite{riscv-spec}.

O registrador \texttt{x0} possui comportamento especial, pois mantém sempre o valor zero, descartando qualquer tentativa de escrita. Os demais registradores podem armazenar operandos, resultados intermediários, endereços e valores utilizados durante a execução do programa. Além da identificação numérica, os registradores também possuem nomes definidos pela ABI, como \texttt{zero}, \texttt{ra}, \texttt{sp}, \texttt{t0}, \texttt{s0} e \texttt{a0}, que facilitam a leitura do código Assembly. Durante a montagem, esses nomes são convertidos para seus respectivos códigos de registrador.

As instruções da arquitetura RISC-V32I são organizadas em formatos padronizados, como R, I, S, B, U e J. Cada formato define a posição dos campos dentro da palavra de 32 bits, incluindo \texttt{opcode}, \texttt{rd}, \texttt{rs1}, \texttt{rs2}, \texttt{funct3}, \texttt{funct7} e campos de imediato. O formato R é usado em operações entre registradores, como \texttt{add} e \texttt{sub}; o formato I em instruções com imediatos, cargas e saltos indiretos; o formato S em armazenamentos na memória; o formato B em desvios condicionais; o formato U em imediatos superiores; e o formato J em saltos incondicionais.

\subsection{Processo de Montagem}

O assembler é responsável por converter instruções escritas em Assembly para código de máquina. Para isso, o programa deve ler o arquivo de entrada, remover comentários e espaços desnecessários, identificar diretivas, rótulos, instruções e dados, além de interpretar registradores, imediatos e operandos.

Quando há rótulos, o assembler também precisa associá-los aos seus respectivos endereços. Essa etapa é importante porque instruções de desvio e salto dependem do cálculo de deslocamentos relativos entre a instrução atual e o destino indicado pelo rótulo. Ao final do processo, os campos de cada instrução são posicionados conforme o formato correspondente da arquitetura RISC-V32I.

\subsection{Seções \texttt{.text} e \texttt{.data}}

Em programas Assembly, as seções \texttt{.text} e \texttt{.data} organizam o conteúdo do arquivo. A seção \texttt{.text} armazena as instruções que serão executadas pelo processador, enquanto a seção \texttt{.data} contém dados estáticos, como constantes, vetores e valores inicializados.

Essa separação é importante porque instruções e dados costumam ocupar regiões distintas da memória. No assembler desenvolvido, o conteúdo da seção \texttt{.text} é convertido para a memória de instruções, enquanto os valores da seção \texttt{.data} são direcionados para a memória de dados.

\subsection{Formato MIF}

O formato MIF, ou \textit{Memory Initialization File}, é utilizado para representar o conteúdo inicial de uma memória em ferramentas de simulação e projetos digitais \cite{mif}. Esse arquivo permite definir previamente os valores armazenados em determinados endereços de memória.

No contexto deste projeto, são gerados arquivos MIF separados para a memória de instruções e para a memória de dados. A memória de instruções recebe os valores hexadecimais correspondentes às instruções codificadas, enquanto a memória de dados recebe os valores declarados na seção \texttt{.data}.

\begin{lstlisting}[style=asmstyle, caption={Estrutura geral de um arquivo MIF.}, label={lst:exemplomif}]
WIDTH=32;
DEPTH=256;

ADDRESS_RADIX=HEX;
DATA_RADIX=HEX;

CONTENT BEGIN
    00 : 00000013;
    01 : 00100093;
END;
\end{lstlisting}

\section{Materiais e Métodos}

\subsection{Materiais Utilizados}

O projeto foi desenvolvido em linguagem Python 3.7 ou superior, em um único arquivo-fonte denominado \texttt{main.py}. Foram utilizadas apenas bibliotecas da própria linguagem, como \texttt{pathlib}, para manipulação de caminhos, \texttt{re}, para reconhecimento de padrões textuais, e \texttt{sys}, para leitura de argumentos pela linha de comando.

Como entrada, a aplicação utiliza arquivos Assembly RISC-V32I com extensão \texttt{.asm}. Como saída, são gerados arquivos no formato \texttt{.mif}, utilizados para inicialização de memória. O simulador RARS foi utilizado como referência para a sintaxe das instruções e para o mapa de memória adotado \cite{rars}, considerando os endereços-base \texttt{TEXT\_BASE = 0x00400000} e \texttt{DATA\_BASE = 0x10010000}.

\subsection{Procedimento Metodológico}

A metodologia adotada consistiu em dividir o processo de montagem em etapas sequenciais. Inicialmente, o arquivo Assembly de entrada é validado, verificando-se a existência do arquivo e a extensão \texttt{.asm}. Em seguida, o conteúdo é lido e submetido a um pré-processamento, no qual são removidos comentários, espaços desnecessários e linhas vazias, preservando caracteres especiais quando presentes dentro de strings.

Após o tratamento inicial, o programa separa o conteúdo nas seções \texttt{.text} e \texttt{.data}. A seção \texttt{.text} é destinada à montagem das instruções, enquanto a seção \texttt{.data} é utilizada para a geração da memória de dados. Para representar o mapa de memória utilizado, foram adotados os endereços-base \texttt{TEXT\_BASE = 0x00400000} e \texttt{DATA\_BASE = 0x10010000}.

A montagem da seção \texttt{.text} foi realizada em duas passagens. Na primeira, são identificados os rótulos e seus respectivos endereços. Na segunda, as instruções são codificadas, utilizando as informações previamente armazenadas para calcular deslocamentos em instruções de desvio e salto. Já a seção \texttt{.data} é processada separadamente, interpretando diretivas como \texttt{.word}, \texttt{.half}, \texttt{.byte}, \texttt{.string}, \texttt{.ascii}, \texttt{.space}, \texttt{.zero} e \texttt{.align}.

\subsection{Validação da Saída}

Ao final do processo, a aplicação gera dois arquivos no formato MIF: um correspondente à memória de instruções, com sufixo \texttt{.text.mif}, e outro correspondente à memória de dados, com sufixo \texttt{.data.mif}. Cada arquivo contém cabeçalho com profundidade, largura da palavra, base dos endereços e base dos dados, utilizando representação hexadecimal.

Além da geração dos arquivos, o programa realiza verificações para identificar erros durante a montagem, como arquivo inexistente, extensão inválida, registrador incorreto, imediato mal formado, operando de memória inválido, rótulo não definido, opcode desconhecido, quantidade inadequada de operandos e diretiva inválida na seção \texttt{.data}. Essas verificações auxiliam na correção do arquivo de entrada e tornam o processo de montagem mais seguro.

\section{Implementação do Assembler}

Esta seção apresenta a estrutura interna da aplicação desenvolvida para realizar a montagem de programas escritos em Assembly RISC-V32I. A implementação foi concentrada em um único arquivo Python, denominado \texttt{main.py}, conforme exigido na especificação do laboratório. O programa foi organizado em etapas responsáveis pela leitura da entrada, tratamento das linhas, separação das seções, montagem das instruções e geração dos arquivos de saída no formato MIF.

\subsection{Estrutura Interna do Código}

A implementação do assembler foi organizada em funções e classes responsáveis por etapas específicas da montagem. As funções iniciais realizam a validação do arquivo de entrada, a leitura do conteúdo, o tratamento das linhas e a separação das seções \texttt{.text} e \texttt{.data}.

Em seguida, o programa utiliza tabelas internas para reconhecer registradores e instruções suportadas. A tabela \texttt{REG\_NAMES} associa os nomes numéricos e os nomes da ABI aos códigos dos registradores, enquanto a tabela \texttt{TABELA\_INSTRUCOES} armazena, para cada mnemônico, o formato da instrução, o \texttt{opcode}, o \texttt{funct3} e, quando necessário, o \texttt{funct7}.

A partir dessas estruturas, a implementação realiza a montagem das instruções e dos dados, gerando ao final os arquivos \texttt{.text.mif} e \texttt{.data.mif}. Dessa forma, a seção de implementação passa a enfatizar a organização interna do código, enquanto a seção de materiais e métodos fica responsável por descrever o procedimento adotado no projeto.

\subsection{Tabelas de Registradores e Instruções}

Para reconhecer os registradores da arquitetura RISC-V32I, foi criada a tabela \texttt{REG\_NAMES}. Essa tabela associa os nomes numéricos dos registradores, como \texttt{x0}, \texttt{x1} e \texttt{x2}, aos nomes convencionais da ABI, como \texttt{zero}, \texttt{ra}, \texttt{sp}, \texttt{t0}, \texttt{s0} e \texttt{a0}. Durante a montagem, cada registrador é convertido para seu respectivo número inteiro entre 0 e 31.

As instruções suportadas foram organizadas na tabela \texttt{TABELA\_INSTRUCOES}. Para cada mnemônico, a tabela armazena o formato da instrução, o \texttt{opcode}, o \texttt{funct3} e, quando necessário, o \texttt{funct7}. Essa organização centraliza as informações de codificação e facilita a seleção da rotina adequada para cada tipo de instrução.

A tabela contempla instruções dos formatos R, I, S, B, U e J, incluindo \texttt{add}, \texttt{sub}, \texttt{and}, \texttt{or}, \texttt{xor}, \texttt{sll}, \texttt{srl}, \texttt{slt}, \texttt{addi}, \texttt{andi}, \texttt{ori}, \texttt{xori}, \texttt{slti}, \texttt{lw}, \texttt{lhu}, \texttt{jalr}, \texttt{sw}, \texttt{beq}, \texttt{bne}, \texttt{lui}, \texttt{auipc} e \texttt{jal}.

\subsection{Codificação das Instruções}

A codificação das instruções foi implementada por meio de classes específicas para os formatos \texttt{TipoR}, \texttt{TipoI}, \texttt{TipoS}, \texttt{TipoB}, \texttt{TipoU} e \texttt{TipoJ}. Cada classe recebe os campos necessários da instrução e os organiza na ordem definida pela arquitetura RISC-V32I, formando uma palavra binária de 32 bits.

A função principal dessa etapa é \texttt{codificar\_instrucao}. Ela recebe o mnemônico, os operandos, o contador de programa atual e o dicionário de rótulos. A partir do mnemônico, o programa consulta a tabela de instruções, identifica o formato correspondente, valida a quantidade de operandos e converte registradores e imediatos.

Nos formatos B e J, utilizados em desvios e saltos, o programa calcula o deslocamento relativo entre a instrução atual e o rótulo de destino. Esse deslocamento é então reorganizado conforme a posição dos bits exigida pelo formato da instrução. Após a montagem dos campos, a instrução é convertida para inteiro e posteriormente escrita em hexadecimal no arquivo MIF.

\subsection{Montagem das Seções \texttt{.text} e \texttt{.data}}

A montagem da seção \texttt{.text} foi realizada em duas passagens. Na primeira passagem, o programa percorre as linhas da seção de instruções e identifica os rótulos, associando cada um ao endereço correspondente. O contador de programa é iniciado em \texttt{TEXT\_BASE = 0x00400000} e incrementado em 4 bytes a cada instrução, pois as instruções RV32I possuem tamanho fixo de 32 bits.

Na segunda passagem, as instruções são efetivamente codificadas. Para cada linha, o mnemônico é separado dos operandos, os registradores e imediatos são interpretados e a função de codificação é chamada. Quando uma instrução utiliza um rótulo como destino, o endereço previamente armazenado é usado para calcular o deslocamento correto.

A seção \texttt{.data} é processada separadamente. Nessa etapa, o programa identifica rótulos e associa cada um a um endereço calculado a partir de \texttt{DATA\_BASE = 0x10010000}. Os dados são armazenados inicialmente em um \texttt{bytearray}, permitindo a inserção de valores byte a byte e o tratamento de diretivas de diferentes tamanhos.

Foram implementadas diretivas como \texttt{.word}, \texttt{.half}, \texttt{.byte}, \texttt{.string}, \texttt{.ascii}, \texttt{.space}, \texttt{.zero} e \texttt{.align}. Ao final do processamento, os bytes são completados até múltiplos de quatro e convertidos em palavras de 32 bits no formato \textit{little-endian}.

\subsection{Geração dos Arquivos MIF}

Após a montagem das seções, o programa gera dois arquivos de saída no formato MIF. O primeiro arquivo, com sufixo \texttt{.text.mif}, armazena as instruções codificadas. O segundo, com sufixo \texttt{.data.mif}, armazena os dados declarados na seção \texttt{.data}.

A função responsável pela geração do arquivo escreve o cabeçalho com a profundidade da memória, a largura da palavra, a base de endereços e a base dos dados. Em seguida, cada palavra é gravada em hexadecimal, associada a um endereço sequencial.

\begin{lstlisting}[style=pythonstyle, caption={Fluxo principal simplificado da aplicação.}, label={lst:main_fluxo}]
conteudo_original = ler_arquivo(caminho_asm)
conteudo_limpo = tratar_ws(conteudo_original)

linhas_text, linhas_data = particionar_secoes(conteudo_limpo)
palavras_data, simbolos_data = montar_data(linhas_data)

simbolos_text = primeira_passagem(linhas_text)
dicionario_labels = {**simbolos_data, **simbolos_text}
palavras_text = segunda_passagem(linhas_text, dicionario_labels)

gerar_arquivo_mif(saida_text, palavras_text, depth=16384)
gerar_arquivo_mif(saida_data, palavras_data, depth=32768)
\end{lstlisting}

\subsection{Tratamento de Erros}

A implementação possui verificações para identificar problemas durante a montagem. Entre os erros tratados estão arquivo inexistente, extensão inválida, registrador inválido, imediato mal formado, operando de memória incorreto, rótulo não definido, opcode desconhecido, quantidade incorreta de operandos e diretiva inválida na seção \texttt{.data}.

Quando um erro é encontrado, o programa interrompe a execução e exibe uma mensagem indicando a causa do problema. Nos casos associados a uma linha específica do arquivo Assembly, a mensagem também informa o número da linha, facilitando a correção pelo usuário.

Portanto, a implementação segue a estrutura básica de um assembler de duas passagens, com tratamento separado para dados e instruções. Essa abordagem permite reconhecer rótulos, calcular deslocamentos relativos, converter instruções para palavras de 32 bits e gerar arquivos compatíveis com o formato MIF.

\section{Resultados e Discussão}

\subsection{Resultados Obtidos}

Como resultado da implementação, foi desenvolvido um assembler em Python capaz de receber um arquivo Assembly com extensão \texttt{.asm}, processar separadamente as seções \texttt{.text} e \texttt{.data} e gerar arquivos de saída no formato MIF. A aplicação utiliza interface por terminal, recebendo o caminho do arquivo de entrada pela linha de comando ou, na ausência desse argumento, utilizando um arquivo padrão.

A seção \texttt{.text} é responsável pela montagem das instruções do programa. Para isso, o assembler identifica rótulos, interpreta mnemônicos, operandos, registradores e imediatos, e converte cada instrução para uma palavra de 32 bits. Ao final, essas palavras são gravadas no arquivo \texttt{.text.mif}.

A seção \texttt{.data}, por sua vez, é processada para armazenar os dados declarados no programa Assembly. Foram tratadas diretivas como \texttt{.word}, \texttt{.half}, \texttt{.byte}, \texttt{.string}, \texttt{.ascii}, \texttt{.space}, \texttt{.zero} e \texttt{.align}. Os dados são organizados em bytes, ajustados quando necessário por alinhamento e convertidos para palavras de 32 bits em formato \textit{little-endian}, sendo gravados no arquivo \texttt{.data.mif}.

\begin{table}[!ht]
\centering
\caption{Resumo dos principais resultados da implementação}
\label{tab:resultados_implementacao}
\begin{tabular}{ll}
\toprule
Aspecto analisado & Resultado obtido \\
\midrule
Linguagem utilizada & Python 3.7 ou superior \\
Arquivo-fonte & Arquivo único \texttt{main.py} \\
Entrada do programa & Arquivo Assembly com extensão \texttt{.asm} \\
Saídas geradas & Arquivos \texttt{.text.mif} e \texttt{.data.mif} \\
Registradores reconhecidos & 32 registradores inteiros e nomes ABI \\
Formatos suportados & R, I, S, B, U e J \\
Seções tratadas & \texttt{.text} e \texttt{.data} \\
Arquivos MIF gerados & Um para instruções e outro para dados \\
\bottomrule
\end{tabular}
\end{table}

\subsection{Instruções e Diretivas Suportadas}

A implementação contemplou as principais instruções solicitadas no laboratório, incluindo operações aritméticas, lógicas, acessos à memória, desvios condicionais e saltos. As instruções foram organizadas em uma tabela interna contendo o formato, o \texttt{opcode}, o \texttt{funct3} e, quando necessário, o \texttt{funct7}.

Entre as instruções implementadas estão \texttt{add}, \texttt{sub}, \texttt{and}, \texttt{or}, \texttt{xor}, \texttt{sll}, \texttt{srl}, \texttt{slt}, \texttt{addi}, \texttt{andi}, \texttt{ori}, \texttt{xori}, \texttt{slti}, \texttt{lw}, \texttt{lhu}, \texttt{sw}, \texttt{beq}, \texttt{bne}, \texttt{jal}, \texttt{jalr}, \texttt{lui} e \texttt{auipc}. Além disso, a implementação também inclui instruções de deslocamento imediato, como \texttt{slli}, \texttt{srli} e \texttt{srai}.

\begin{table}[!ht]
\centering
\caption{Síntese quantitativa das funcionalidades implementadas}
\label{tab:sintese_quantitativa}
\begin{tabular}{ll}
\toprule
Métrica & Valor observado \\
\midrule
Registradores inteiros reconhecidos & 32 \\
Formatos de instrução tratados & 6 \\
Mnemônicos presentes na tabela interna & 25 \\
Diretivas principais da seção \texttt{.data} & 8 \\
Arquivos de saída gerados & 2 \\
Passagens realizadas na seção \texttt{.text} & 2 \\
\bottomrule
\end{tabular}
\end{table}

Esses resultados indicam que o assembler desenvolvido atende ao objetivo central do laboratório, pois realiza a conversão das instruções Assembly para representação hexadecimal em arquivos MIF e também trata dados declarados na seção \texttt{.data}.

\subsection{Avaliação do Funcionamento}

O funcionamento do assembler foi organizado em etapas bem definidas. Primeiro, o arquivo de entrada é validado e lido. Em seguida, comentários e espaços desnecessários são removidos, preservando caracteres especiais dentro de strings. Depois disso, o conteúdo é separado entre as seções \texttt{.text} e \texttt{.data}.

Na montagem da seção \texttt{.text}, o uso de duas passagens foi um resultado importante da implementação. A primeira passagem permite identificar os rótulos e associá-los aos seus endereços. A segunda passagem realiza a codificação das instruções, já com os endereços dos rótulos disponíveis. Essa abordagem permite tratar corretamente instruções de desvio e salto que referenciam rótulos definidos antes ou depois da instrução atual.

Na seção \texttt{.data}, o uso de um \texttt{bytearray} permitiu armazenar dados de diferentes tamanhos antes da conversão final para palavras de 32 bits. Isso facilitou o tratamento de diretivas como \texttt{.byte}, \texttt{.half}, \texttt{.word} e strings, além de permitir a aplicação de alinhamento quando necessário.

\subsection{Discussão dos Resultados}

Os resultados obtidos mostram que a implementação conseguiu reproduzir as etapas fundamentais de um assembler. O programa realiza a leitura do código Assembly, interpreta registradores e imediatos, reconhece rótulos, codifica instruções em palavras de 32 bits e gera arquivos compatíveis com o formato MIF.

Um ponto positivo da implementação é a separação clara entre as etapas do processo de montagem. A existência de funções específicas para leitura, tratamento das linhas, particionamento das seções, montagem de dados, identificação de rótulos, codificação de instruções e geração dos arquivos MIF torna o código mais organizado e facilita futuras correções ou expansões.

Outro aspecto relevante é o tratamento de erros. A aplicação identifica problemas como arquivo inexistente, extensão inválida, registrador incorreto, imediato mal formado, operando de memória inválido, opcode desconhecido, quantidade inadequada de operandos e rótulos não definidos. Isso evita que entradas inválidas sejam processadas sem aviso e facilita a localização de falhas no arquivo Assembly.

Em relação ao desempenho, a implementação apresenta comportamento adequado para o escopo do laboratório. A seção \texttt{.data} é percorrida uma vez, enquanto a seção \texttt{.text} é percorrida duas vezes. Dessa forma, o tempo de processamento cresce de maneira aproximadamente linear em relação à quantidade de linhas do arquivo de entrada, o que é apropriado para arquivos Assembly de pequeno e médio porte.

Como pontos de melhoria, seria possível ampliar a quantidade de testes automatizados, comparar os arquivos MIF gerados com saídas de referência e padronizar melhor as mensagens exibidas ao usuário. Também seria possível aperfeiçoar a geração dos arquivos MIF para preencher de forma mais completa os endereços não utilizados ou utilizar uma representação mais compacta para regiões de memória preenchidas com zero.


\section{Conclusão}

O desenvolvimento do assembler permitiu aplicar, de forma prática, conceitos fundamentais da arquitetura RISC-V32I, como codificação de instruções, uso de registradores, organização de memória, tratamento de rótulos e geração de arquivos de inicialização no formato MIF. A implementação em Python possibilitou transformar arquivos Assembly em representações hexadecimais compatíveis com memórias de instruções e de dados.

De modo geral, o projeto atingiu seu objetivo principal ao realizar a leitura de arquivos \texttt{.asm}, processar as seções \texttt{.text} e \texttt{.data}, codificar instruções dos principais formatos da arquitetura e gerar arquivos \texttt{.mif} separados. Além disso, o uso de duas passagens na seção de instruções mostrou-se essencial para resolver rótulos e calcular corretamente deslocamentos em desvios e saltos.

Apesar dos resultados satisfatórios, a implementação ainda pode ser aprimorada com a inclusão de mais testes automatizados, melhorias nas mensagens de erro e comparação direta com saídas de referência. Ainda assim, o trabalho cumpriu sua finalidade ao demonstrar o funcionamento básico de um assembler e reforçar a compreensão da relação entre código Assembly, código de máquina e organização da memória.


\begin{thebibliography}{00}

\bibitem{riscv-spec}
RISC-V International, \textit{The RISC-V Instruction Set Manual}. [Online]. Available: \url{https://riscv.org/technical/specifications/}

\bibitem{rars}
The Third One, \textit{RARS: RISC-V Assembler and Runtime Simulator}. [Online]. Available: \url{https://github.com/TheThirdOne/rars}

\bibitem{mif}
Intel, \textit{Memory Initialization File (.mif) Definition}. [Online]. Available: \url{https://www.intel.com/content/www/us/en/programmable/quartushelp/17.0/reference/glossary/def_mif.htm}

\end{thebibliography}

\end{document}
